% ID: SEL-D0021-OB-E03
% ESTADO: validada. S022, revisada con los criterios oficiales.
Sean $T$ y $P$ aprobar los exámenes teórico y práctico.
\textbf{a)} $P(T\cup P)=0.8+0.6-0.5=\boxed{0.9}$, restando la
intersección para contarla una sola vez.

\textbf{b)} De los que aprueban el práctico, $0.6-0.5=0.1$ de
la población no aprueba el teórico. Dividimos por $P(T^c)=0.2$:
$P(P\mid T^c)=\boxed{0.5}$.

\textbf{c)} La independencia exigiría $P(T\cap P)=P(T)P(P)$.
Como $0.5\ne0.8\cdot0.6=0.48$, los sucesos son dependientes.
